\documentclass[12pt]{article}
\usepackage[utf8]{inputenc}
\usepackage[T1]{fontenc}
\usepackage[margin=1in]{geometry}
\usepackage{setspace}
\usepackage{xcolor}
\usepackage[colorlinks=true, urlcolor=blue!60!black, linkcolor=blue!60!black]{hyperref}
\setlength{\parindent}{2em}
\setlength{\parskip}{1.25em}
\renewcommand{\baselinestretch}{1.0}
\title{\textit{Gundy v. United States}, 588 U.S. 128 (2019)}
\author{Case Brief}
\date{}

\begin{document}
\maketitle

\section{Case Metadata}
\begin{description}
  \item[Full Case Name:] \textit{Herman Avery Gundy v. United States}
  \item[Citation:] 588 U.S. 128 (2019) (Docket No. 17-6086)
  \item[Argued / Decided:] October 2, 2018 / June 20, 2019
  \item[Vote:] 5--3 (4--1--3 plurality judgment; Justice Kavanaugh took no part in the consideration or decision of the case)
  \item[Author:] Justice Elena Kagan (plurality opinion announcing the judgment of the Court, joined by Justices Ginsburg, Breyer, and Sotomayor)
  \item[Concurrences / Dissents:] Justice Alito (concurring in the judgment); Justice Gorsuch (dissenting, joined by Chief Justice Roberts and Justice Thomas); Justice Kavanaugh (not participating)
  \item[Chapter / Topic:] Separation of Powers and the Executive --- The Nondelegation Doctrine and Article I Legislative Power
  \item[Primary Source URL:] \url{https://supreme.justia.com/cases/federal/us/588/17-6086/}
\end{description}

\section{Facts and Procedural History}
In 2006, Congress enacted the Sex Offender Registration and Notification Act (SORNA), 120 Stat. 590, 34 U.S.C. \S~20901 \textit{et seq.}, to replace an inconsistent patchwork of state sex-offender registries---under which over 100,000 sex offenders had escaped registration---with ``a comprehensive national system for the registration of those offenders,'' backed by federal felony penalties of up to ten years' imprisonment for knowingly failing to register after interstate travel (18 U.S.C. \S~2250(a)). While 34 U.S.C. \S~20913(b) established the general rule that offenders must register before completing their prison sentences (or within three business days of sentencing if not imprisoned), Congress confronted a transitional dilemma for the roughly 500,000 ``pre-Act offenders'' convicted before SORNA's enactment, many of whom had already been released from custody and could not comply with \S~20913(b).

To address pre-Act offenders, Congress enacted 34 U.S.C. \S~20913(d) (formerly 42 U.S.C. \S~16913(d)), titled ``Initial registration of sex offenders unable to comply with subsection (b),'' which provided: ``The Attorney General shall have the authority to specify the applicability of the requirements of this subchapter to sex offenders convicted before the enactment of this chapter \dots\ and to prescribe rules for the registration of any such sex offenders and for other categories of sex offenders who are unable to comply with subsection (b).'' In February 2007 (217 days after enactment), the Attorney General issued an interim rule applying SORNA's registration requirements in full to all pre-Act offenders (72 Fed. Reg. 8897), and finalized that rule in December 2010 (75 Fed. Reg. 81850).

In 2005---the year before SORNA's enactment---petitioner Herman Avery Gundy pleaded guilty in Maryland to sexually assaulting a minor. Following his release from custody in 2012, Gundy traveled to New York and resided there without registering as a sex offender. He was indicted and convicted in the U.S. District Court for the Southern District of New York under 18 U.S.C. \S~2250. Gundy moved to dismiss, arguing that \S~20913(d) violated Article I's nondelegation doctrine by granting the Attorney General unguided, plenary authority to decide whether and how SORNA applied to pre-Act offenders. Both the District Court and the U.S. Court of Appeals for the Second Circuit rejected his nondelegation claim, 695 Fed. Appx. 639 (2017), joining all eleven Courts of Appeals to have considered the issue. The Supreme Court granted certiorari limited to the nondelegation question, 583 U.S. 1166 (2018). Because Justice Kavanaugh was confirmed after the October 2, 2018 oral argument, an eight-member Court decided the case.

\section{Issue(s) Presented}
Does 34 U.S.C. \S~20913(d) of the Sex Offender Registration and Notification Act violate Article I, Section 1's nondelegation doctrine by authorizing the Attorney General to ``specify the applicability'' of SORNA's registration requirements to sex offenders convicted prior to the Act's enactment?

\section{Holding}
No. In a 5--3 judgment (with no majority opinion), the Supreme Court affirmed Gundy's conviction, holding that \S~20913(d) does not violate the nondelegation doctrine. A four-Justice plurality construed \S~20913(d) to require the Attorney General to apply SORNA's registration requirements to all pre-Act offenders ``as soon as feasible,'' a narrow transitional delegation that easily satisfied the ``intelligible principle'' test. Justice Alito supplied the decisive fifth vote concurring solely in the judgment, concluding that the statute did not fail the Court's deferential post-1935 nondelegation precedent, while signaling his willingness to reconsider that doctrine in a future case with a full Court.

\section{Reasoning}
Writing for the four-Justice plurality, Justice Kagan began with the constitutional baseline: Article I, \S~1 vests ``[a]ll legislative Powers herein granted'' in Congress, barring Congress from transferring powers that are ``strictly and exclusively legislative'' (\textit{Wayman v. Southard}, 10 Wheat. 1, 42--43 (1825)). Yet, recognizing that ``Congress simply cannot do its job absent an ability to delegate power under broad general directives'' (\textit{Mistretta v. United States}, 488 U.S. 361, 372 (1989)), the Court has long held that a statutory delegation is constitutional so long as Congress lays down an ``intelligible principle'' to guide the executive official's discretion (\textit{J.W. Hampton, Jr., \& Co. v. United States}, 276 U.S. 394, 409 (1928)).

\textbf{1. Statutory Interpretation as the Dispositive First Step:} Justice Kagan observed that ``a nondelegation inquiry always begins (and often almost ends) with statutory interpretation.'' Gundy argued that the first half of \S~20913(d) gave the Attorney General ``unguided'' and ``unchecked'' authority to require all, some, or none of the pre-Act offenders to register, or to reverse course at whim. Justice Kagan conceded that ``[i]f that were so, we would face a nondelegation question.'' However, reading \S~20913(d) holistically in light of text, context, purpose, history, and precedent, the plurality rejected Gundy's reading.

\textbf{2. Reliance on \textit{Reynolds v. United States} and Textual Context:} In \textit{Reynolds v. United States}, 565 U.S. 432 (2012), the Court had held that SORNA did not apply to pre-Act offenders until the Attorney General so specified, explaining that Congress had created a temporary stopgap because ``instantaneous registration'' of hundreds of thousands of already-released offenders raised practical implementation problems. Revisiting the statute independently, Justice Kagan pointed to four textual and contextual cues confirming that the Attorney General's discretion extended only to \textit{feasibility} and timing:
\begin{enumerate}
  \item \textit{Declaration of Purpose (\S~20901):} Congress declared its purpose to establish a ``comprehensive national system''---an ``all-encompassing'' mandate incompatible with permanently exempting the majority of existing sex offenders.
  \item \textit{Past-Tense Definition (\S~20911(1)):} Defining a ``sex offender'' as an individual who ``\textit{was} convicted'' of a sex offense confirmed Congress's intent to capture pre-Act offenders.
  \item \textit{Legislative History:} Committee reports and floor statements focused intensely on capturing the more than 100,000 ``missing'' pre-Act sex offenders.
  \item \textit{Title and Second Clause of \S~20913(d):} Addressing offenders ``unable to comply with subsection (b)'' revealed that ``specify the applicability'' meant ``specify \textit{how} to apply SORNA'' during the transition, not ``specify \textit{whether} to apply SORNA.'' Moreover, every Attorney General had applied SORNA to all pre-Act offenders without exception.
\end{enumerate}

\textbf{3. Application of the Intelligible Principle Test:} Once \S~20913(d) was construed as a time-limited directive to register pre-Act offenders ``as soon as feasible,'' the constitutional inquiry was ``easy.'' The Court had invalidated statutes on nondelegation grounds only twice in American history---both in 1935 (\textit{A.L.A. Schechter Poultry Corp. v. United States}, 295 U.S. 495 (1935), and \textit{Panama Refining Co. v. Ryan}, 293 U.S. 388 (1935))---where Congress articulated no policy or standard whatsoever. By contrast, the Court had routinely upheld far broader delegations (``public interest,'' ``just and reasonable,'' ``requisite to protect the public health''). Compared to those precedents, SORNA's transitional feasibility delegation was ``distinctly small-bore,'' and invalidating it would imply that ``most of Government is unconstitutional.''

\section{Concurring and Dissenting Opinions}
\begin{itemize}
  \item \textbf{Justice Alito (Concurring in the Judgment):} Issued a terse, pivotal three-paragraph opinion declining to join either the plurality's statutory interpretation or its constitutional defense of broad delegation. Observing that since 1935 the Court had uniformly upheld ``extraordinarily capacious standards,'' Justice Alito stated that while he would support a majority willing to reconsider the Court's 84-year-old nondelegation approach, ``it would be freakish to single out the provision at issue here for special treatment'' absent five votes to overhaul the doctrine.
  \item \textbf{Justice Gorsuch (Dissenting, joined by Chief Justice Roberts and Justice Thomas):} Delivered a sweeping originalist manifesto calling to replace the modern ``intelligible principle'' test with the Framers' tripartite test for Article I legislative power. Grounding nondelegation in John Locke, Montesquieu, bicameralism, presentment, fair notice, and democratic accountability, Justice Gorsuch argued that Congress may delegate authority to the Executive in only three constitutionally permissible categories:
  \begin{enumerate}
    \item \textit{``Fill Up the Details'':} Once Congress makes the core policy decisions governing private conduct, it may authorize the Executive to fill up administrative details (\textit{Wayman v. Southard}).
    \item \textit{Conditional Fact-Finding:} Congress may prescribe a rule governing private conduct whose application depends on executive fact-finding (\textit{Cargo of the Brig Aurora v. United States}, 7 Cranch 382 (1813)).
    \item \textit{Overlapping Non-Legislative Responsibilities:} Congress may confer broad discretion over matters already within the Executive's or Judiciary's independent constitutional powers (e.g., foreign affairs or court rules).
  \end{enumerate}
  Justice Gorsuch argued that Chief Justice Taft's 1928 ``intelligible principle'' phrase in \textit{J.W. Hampton} was merely a shorthand summary of these traditional limits that later ``mutated'' into a toothless rubber stamp, forcing the Court to police delegations indirectly through the major questions doctrine and void-for-vagueness cases. Turning to SORNA, he accused the plurality of rewriting \S~20913(d) to invent a ``feasibility'' standard absent from the text---contradicting the government's own prior representations in \textit{Reynolds} that the Attorney General had unfettered discretion---and warned against uniting lawmaking and prosecutorial power in the nation's chief prosecutor.
\end{itemize}

\section{Key Quotations}
\begin{quote}
``Given that standard, a nondelegation inquiry always begins (and often almost ends) with statutory interpretation. Only after a court has determined a challenged statute's meaning can it decide whether the law sufficiently guides executive discretion to accord with Article I.'' \\
--- Justice Kagan, plurality opinion, 588 U.S. at 135--136.
\end{quote}

\begin{quote}
``Indeed, if SORNA's delegation is unconstitutional, then most of Government is unconstitutional---dependent as Congress is on the need to give discretion to executive officials to implement its programs.'' \\
--- Justice Kagan, plurality opinion, 588 U.S. at 149.
\end{quote}

\begin{quote}
``If a majority of this Court were willing to reconsider the approach we have taken for the past 84 years, I would support that effort. But because a majority is not willing to do that, it would be freakish to single out the provision at issue here for special treatment.'' \\
--- Justice Alito, concurring in the judgment, 588 U.S. at 149.
\end{quote}

\begin{quote}
``The Constitution promises that only the people's elected representatives may adopt new federal laws restricting liberty. Yet the statute before us scrambles that design. It purports to endow the nation's chief prosecutor with the power to write his own criminal code governing the lives of a half-million citizens.'' \\
--- Justice Gorsuch, dissenting, 588 U.S. at 149--150.
\end{quote}

\begin{quote}
``Although it is nominally a canon of statutory construction, we apply the major questions doctrine in service of the constitutional rule that Congress may not divest itself of its legislative power by transferring that power to an executive agency.'' \\
--- Justice Gorsuch, dissenting, 588 U.S. at 167.
\end{quote}

\section{Doctrinal Significance and Modern Legacy (Through 2026)}
Although the government won the judgment in \textit{Gundy}, the decision initially appeared to put the modern nondelegation doctrine on life support: four Justices (Gorsuch, Roberts, Thomas, and Alito) expressed readiness to reconsider the post-1935 intelligible-principle framework, and five months later Justice Kavanaugh---who had been recused in \textit{Gundy}---issued a statement respecting the denial of certiorari in \textit{Paul v. United States}, 589 U.S. \_\_\_ (Nov. 25, 2019), praising Justice Gorsuch's ``scholarly analysis'' as warranting ``further consideration in future cases'' (\url{https://www.supremecourt.gov/opinions/19pdf/17-8830_5hdk.pdf}).

Yet between 2020 and 2026, the predicted constitutional revolution took a markedly different trajectory---vindicating Justice Gorsuch's ``hydraulic pressures'' observation while preserving the intelligible-principle test on the books:
\begin{itemize}
  \item \textbf{Indirect Policing via Statutory Canons (2022--2024):} Rather than invalidating statutes under Article I, the Court aggressively constrained agency lawmaking by elevating the Major Questions Doctrine in \textit{West Virginia v. EPA}, 597 U.S. 697 (2022), and \textit{Biden v. Nebraska}, 600 U.S. 477 (2023), and by overruling \textit{Chevron} deference in \textit{Loper Bright Enterprises v. Raimondo}, 603 U.S. 369 (2024). (In \textit{Allstates Refractory Contractors, LLC v. Su}, July 2, 2024, Justice Thomas dissented from denial of certiorari, unsuccessfully urging the Court to use OSHA's delegation to revisit \textit{Gundy}; \url{https://www.supremecourt.gov/opinions/23pdf/23-819_m648.pdf}.)
  \item \textbf{Reaffirmation of the Intelligible Principle Test in \textit{FCC v. Consumers' Research}, 606 U.S. \_\_\_ (2025):} On June 27, 2025, the Court confronted its first full-dress nondelegation merits case since \textit{Gundy} and decisively rejected the anticipated revival, upholding the FCC's Universal Service Fund (USF) contribution scheme under 47 U.S.C. \S~254 by a 6--3 vote (\url{https://supreme.justia.com/cases/federal/us/606/24-354/}). Justice Kagan---joined by Chief Justice Roberts and Justices Sotomayor, Kavanaugh, Barrett, and Jackson---reaffirmed the ``intelligible principle'' test, rejected a stricter numeric-cap requirement for revenue-raising delegations (following \textit{J.W. Hampton} and \textit{Skinner v. Mid-America Pipeline Co.}, 490 U.S. 212 (1989)), held that ``sufficient'' funding coupled with \S~254's mandatory service criteria constrained the FCC, rejected the private nondelegation challenge to the Universal Service Administrative Company under \textit{Sunshine Anthracite Coal Co. v. Adkins}, 310 U.S. 381 (1940), and rejected the en banc Fifth Circuit's ``combination'' theory. Crucially, Justice Kavanaugh's concurrence explained why the intelligible-principle test endures: \textit{Loper Bright} and the major questions doctrine already check domestic agency overreach, while broad delegations remain vital in foreign affairs and national security; the real accountability threat, he argued, comes from delegations to \textit{independent} agencies insulated from presidential removal. Justice Gorsuch (joined by Justices Thomas and Alito) dissented, repeating his \textit{Gundy} critique that the majority rewrote \S~254's text to save an unprecedented delegation of the taxing power (while noting the majority's footnote 9 refusal to bless \S\S~254(c)(3) and (h)(2)).
  \item \textbf{2026 Consolidation --- Statutory Restraint and Unitary Executive Accountability:} Two blockbuster 2026 decisions cemented the post-\textit{Gundy} equilibrium. First, in \textit{Learning Resources, Inc. v. Trump}, 607 U.S. \_\_\_ (Feb. 20, 2026), the Court struck down the President's IEEPA tariffs 6--3 on statutory grounds---with Chief Justice Roberts and Justices Gorsuch and Barrett invoking the major questions doctrine---rather than relying on constitutional nondelegation (\url{https://supreme.justia.com/cases/federal/us/607/24-1287/}). Second, in \textit{Trump v. Slaughter}, 609 U.S. \_\_\_ (June 29, 2026), the Court overruled \textit{Humphrey's Executor v. United States}, 295 U.S. 602 (1935) and invalidated statutory for-cause removal protections for FTC commissioners (\url{https://supreme.justia.com/cases/federal/us/609/25-332/})---adopting the first of Justice Kavanaugh's two Article II remedies from \textit{Consumers' Research} by restoring presidential removal control over agencies rather than dismantling Congress's ability to delegate under the intelligible-principle test.
\end{itemize}

\section{Pedagogical Critiques and Exam Hypotheticals}
\textbf{Pedagogical Critiques:}
\begin{itemize}
  \item \textit{Constitutional Avoidance vs. Textual Fidelity:} Both \textit{Gundy} (reading ``as soon as feasible'' into \S~20913(d)) and \textit{FCC v. Consumers' Research} (reading \S~254(c)(1)'s ``consider'' factors as separately mandatory) illustrate a recurring critique raised by Justice Gorsuch: to avoid striking down statutes under the nondelegation doctrine, the Court aggressively narrows---or effectively rewrites---broad statutory text. Students should debate whether such saving constructions respect Congress or usurp the legislative drafting pen.
  \item \textit{Criminal vs. Civil Delegations:} \textit{Gundy} involved a delegation defining the scope of a 10-year federal felony, yet the plurality applied the exact same lenient intelligible-principle standard used for civil utility ratemaking, leaving unresolved the question reserved in \textit{Touby v. United States}, 500 U.S. 160 (1991), whether criminal delegations demand heightened congressional specificity.
  \item \textit{The Trade-off Between Nondelegation and Removal Power:} Comparing \textit{Gundy}, \textit{FCC v. Consumers' Research} (2025), and \textit{Trump v. Slaughter} (2026) reveals a deep structural irony: because the Court found a workable constitutional nondelegation line elusive, it instead strengthened presidential control over agencies (eliminating \textit{Humphrey's Executor}) and judicial control over statutory meaning (eliminating \textit{Chevron} and expanding the major questions doctrine).
\end{itemize}

\textbf{Exam Hypothetical (and ``Opposite Side'' Contrast):} \\
\textit{Hypothetical:} Responding to rapid developments in autonomous vehicles and artificial intelligence, Congress enacts the ``National AI Safety Act of 2026.'' Section 4 of the Act empowers the Secretary of Commerce ``to promulgate such rules, licensing conditions, and criminal prohibitions (punishable by up to five years' imprisonment) governing the training and deployment of frontier artificial intelligence models as the Secretary determines are necessary and appropriate to protect the public interest, promote American innovation, and safeguard national security.'' The statute contains no further definition of ``frontier'' models, no threshold metrics, and no list of prohibited practices. Relying on Section 4, the Secretary issues a rule making it a federal felony to release any open-weight AI model trained using more than $10^{25}$ floating-point operations without a federal license. A software developer indicted under the rule challenges Section 4 under the nondelegation doctrine.

\textit{Analysis \& ``Opposite Side'' Contrast:}
\begin{itemize}
  \item \textit{Challenger's Argument (Under Gorsuch's \textit{Gundy} Dissent and the Sliding Scale of \textit{Whitman} / \textit{Consumers' Research}):} Unlike SORNA's transitional ``as soon as feasible'' mandate in \textit{Gundy} or the bounded beneficiary and service criteria in \textit{FCC v. Consumers' Research}, Section 4 confers economy-wide criminal lawmaking power with zero determinate boundaries. Even under the majority's intelligible-principle test in \textit{Consumers' Research}, ``the degree of agency discretion that is acceptable varies according to the scope of the power congressionally conferred'' (\textit{Whitman v. American Trucking Assns., Inc.}, 531 U.S. 457, 475 (2001)). Combining the ``public interest'' with competing goals (``innovation'' vs. ``safety'') without statutory definitions replicates the ``preface of generalities'' struck down in \textit{Schechter Poultry}. Moreover, under the Major Questions Doctrine (\textit{West Virginia v. EPA}; \textit{Learning Resources, Inc. v. Trump}), Congress does not hide an economy-reshaping criminal ban on open-source software in vague boilerplate.
  \item \textbf{Opposite Side Contrast (Government's Defense under Kagan's \textit{Gundy} Plurality and Kavanaugh's \textit{Consumers' Research} Concurrence):} The government will invoke the Court's unbroken post-1935 practice of upholding ``public interest'' standards (\textit{NBC v. United States}; \textit{Yakus}) and argue that ``frontier'' AI and ``national security'' supply contextual boundaries---especially because, as Justice Kavanaugh emphasized in \textit{Consumers' Research} (and Justice Gorsuch acknowledged in his third \textit{Gundy} category), nondelegation and the major questions doctrine play a far more limited role when Congress delegates to the President or an accountable Cabinet Secretary in the \textit{national security} realm. Furthermore, unlike an independent agency, the Secretary of Commerce is removable at will by the President, mitigating Article II accountability concerns. Finally, the government will urge the Court to apply \textit{Gundy}'s statutory-interpretation-first method by construing ``national security'' and ``frontier'' narrowly to cover only demonstrable dual-use biological or cyber weapons risks, thereby saving Section 4 from facial invalidation.
\end{itemize}

\end{document}
